% !TeX root = ../sections/03_solutions.tex
\subsection{対策-シュミットの直交化法}
\begin{exercise}[シュミットの直交化法]
	閉区間 \([-1,1]\) 上の関数
	\[
	\psi_1(x)=x,\qquad
	\psi_2(x)=x^2,\qquad
	\psi_3(x)=1
	\]
	を考える。
	
	関数 \(\psi_k(x)\) と \(\psi_l(x)\) の内積を
	\[
	(\psi_k,\psi_l)
	=
	\int_{-1}^{1}\psi_k(x)\psi_l(x)\,dx
	\]
	で定める。
	
	正規化を伴わないシュミットの直交化法により，
	関数列
	\[
	[\psi_1,\psi_2,\psi_3]
	\]
	から直交関数列
	\[
	[\varphi_1,\varphi_2,\varphi_3]
	\]
	を作れ。
\end{exercise}

\begin{background}
	関数の内積とは，二つの関数を掛けて積分したものである。
	この問題では，
	\[
	(f,g)=\int_{-1}^{1}f(x)g(x)\,dx
	\]
	で定義されている。
	
	二つの関数 \(f,g\) が直交するとは，
	\[
	(f,g)=0
	\]
	が成り立つことである。
	
	これはベクトルの内積で
	\[
	\vec{u}\cdot \vec{v}=0
	\]
	となることの関数版である。
	
	\medskip
	
	シュミットの直交化法は，与えられた関数列
	\[
	[\psi_1,\psi_2,\psi_3]
	\]
	から，互いに直交する関数列
	\[
	[\varphi_1,\varphi_2,\varphi_3]
	\]
	を作る方法である。
	
	正規化を伴わない場合は，長さを \(1\) にそろえる必要はない。
	つまり，
	\[
	(\varphi_i,\varphi_j)=0\quad (i\neq j)
	\]
	だけを満たせばよい。
	
	一方，正規化を伴う場合は，
	\[
	\|e_i\|=1
	\]
	も満たすようにする。
	ここで，
	\[
	\|f\|=\sqrt{(f,f)}
	\]
	である。
\end{background}
\begin{strategy}
	正規化を伴わないシュミットの直交化法は，次の形で進める。
	
	まず，
	\[
	\varphi_1=\psi_1
	\]
	とする。
	
	次に，
	\[
	\varphi_2
	=
	\psi_2
	-
	\frac{(\psi_2,\varphi_1)}{(\varphi_1,\varphi_1)}
	\varphi_1
	\]
	とする。
	
	最後に，
	\[
	\varphi_3
	=
	\psi_3
	-
	\frac{(\psi_3,\varphi_1)}{(\varphi_1,\varphi_1)}
	\varphi_1
	-
	\frac{(\psi_3,\varphi_2)}{(\varphi_2,\varphi_2)}
	\varphi_2
	\]
	とする。
	
	この問題では，
	\[
	\varphi_1=x
	\]
	である。
	
	次に，
	\[
	(\psi_2,\varphi_1)
	=
	(x^2,x)
	=
	\int_{-1}^{1}x^3\,dx
	=
	0
	\]
	であるから，
	\[
	\varphi_2=x^2.
	\]
	
	最後に，
	\[
	\varphi_3
	=
	1
	-
	\frac{(1,x)}{(x,x)}x
	-
	\frac{(1,x^2)}{(x^2,x^2)}x^2
	\]
	を計算する。
	
	ここで，
	\[
	(1,x)=\int_{-1}^{1}x\,dx=0,
	\]
	\[
	(1,x^2)=\int_{-1}^{1}x^2\,dx=\frac{2}{3},
	\]
	\[
	(x^2,x^2)=\int_{-1}^{1}x^4\,dx=\frac{2}{5}.
	\]
	
	したがって，
	\[
	\frac{(1,x^2)}{(x^2,x^2)}
	=
	\frac{2/3}{2/5}
	=
	\frac{5}{3}.
	\]
	
	よって，
	\[
	\varphi_3
	=
	1-\frac{5}{3}x^2.
	\]
	
	したがって，正規化を伴わない直交関数列は
	\[
	\boxed{
		\varphi_1=x,\qquad
		\varphi_2=x^2,\qquad
		\varphi_3=1-\frac{5}{3}x^2
	}
	\]
	である。
	
	分数を避けたい場合は，\(\varphi_3\) を定数倍して
	\[
	\boxed{
		\varphi_1=x,\qquad
		\varphi_2=x^2,\qquad
		\varphi_3=3-5x^2
	}
	\]
	としてもよい。
	直交性は定数倍しても保たれる。
\end{strategy}